\documentclass[12pt]{article}
\usepackage[utf8]{inputenc}
\usepackage[T1]{fontenc}
\usepackage{hyperref}
\title{Tumleo Language Resources for K-12 Curriculum Development}
\author{Compiled by Claude Sonnet 5 (claude-sonnet-5) for Andr\'as Kornai}
\date{}
\begin{document}
\maketitle

\section{Introduction}
Tumleo (ISO 639-3 tmq; Austronesian, Schouten branch of North New Guinea) is spoken on Tumleo Island and the nearby Aitape coast of Sandaun Province, Papua New Guinea. Wikipedia lists about 790 speakers (2003). The web search found only reference-style sources. It found no school materials, primer, grammar or dictionary.

\section{Primary Educational Resources}
No primers, textbooks or school programmes for Tumleo turned up in the search.

\section{Linguistic and Reference Materials}
\begin{itemize}
\item Wikipedia, ``Tumleo language'': overview, classification and speaker figures. \url{https://en.wikipedia.org/wiki/Tumleo_language}
\item ASJP, Tumleo wordlist. It credits Don Laycock's papers (PARADISEC archive) and a wordlist compiled by Matthew Dryer. \url{https://asjp.clld.org/languages/TUMLEO_3}
\item Endangered Languages catalogue entry (endangered; low confidence rating). \url{https://endangeredlanguages.com/lang/10521}
\end{itemize}

\section{Religious and Community Resources}
\begin{itemize}
\item Joshua Project profile of the Tumleo people. It is mainly a prayer profile with secondary-source figures. \url{https://joshuaproject.net/index.php/people_groups/15636}
\end{itemize}

\section{Digital and Audio Resources}
No audio or digital resources were confirmed. The ASJP page points to Laycock material in the PARADISEC archive, but the search did not confirm its contents.

\section{Curriculum Development Context}
Very little documented material was found. A curriculum would first need fieldwork or archive work, for example the PARADISEC Laycock collection and Glottolog's bibliography, to build a lexicon and a basic description. This should be done with the community.

\end{document}